\chapter{Onglet Cp / Couche limite}
\label{chap:cp}

L'onglet \menu{Cp / Couche limite} offre une vue détaillée, façon
\emph{XFoil}, d'un \textbf{calcul précis} : un profil, un Reynolds et
une incidence. Il réunit dans une même fenêtre la distribution de
pression (visqueuse \emph{et} non visqueuse) et la géométrie du profil
entourée de sa couche limite.

\screenshot{20_tab_cp.png}{1.0}{%
    Onglet \menu{Cp / Couche limite} pour un NACA~2412 à
    $Re = 10^6$ et $\alpha = 6\textdegree$. En haut : le coefficient
    de pression $C_p$ (extrados en jaune, intrados en cyan) et sa
    valeur non visqueuse (pointillé blanc). En bas : le profil à
    l'échelle, avec les frontières de couche limite $\delta^*$
    extrados (jaune) et intrados (cyan). L'encart résume les
    coefficients au point choisi.}

\section{Sélection du calcul}

Une barre de contrôle, en haut de l'onglet, permet de cibler le
calcul à visualiser :

\begin{tabularx}{\linewidth}{l X}
    \toprule
    \textbf{Liste} & \textbf{Rôle} \\
    \midrule
    \menu{Profil}
        & Profil à afficher : courant (bleu), référence (rouge) ou
          volet (vert). Seuls les profils effectivement simulés sont
          proposés. \\
    \menu{Re}
        & Nombre de Reynolds du calcul. \\
    \menu{$\alpha$}
        & Incidence (en degrés) du calcul. \\
    \bottomrule
\end{tabularx}

Les listes sont \textbf{en cascade} : changer de profil met à jour la
liste des Reynolds disponibles, puis celle des incidences.

\begin{noteinfo}
Contrairement à l'onglet \menu{Résultats} qui superpose les profils,
cet onglet affiche \textbf{un seul calcul à la fois}, pour une lecture
fine de la pression et de la couche limite.
\end{noteinfo}

\section{Tracé supérieur : coefficient de pression}

Le tracé du haut reprend la convention \emph{XFoil} : l'axe vertical
porte le coefficient de pression $C_p$ \textbf{orienté vers le bas}
— la succion ($C_p < 0$) est donc dirigée vers le \emph{haut}. Trois
courbes y sont superposées :

\begin{itemize}
    \item l'\textbf{extrados} (jaune) ;
    \item l'\textbf{intrados} (cyan) ;
    \item le $C_p$ \textbf{non visqueux} (pointillé blanc), issu d'un
        calcul XFoil sans couche limite (méthode des panneaux,
        indépendant du Reynolds).
\end{itemize}

L'écart entre le $C_p$ visqueux et le $C_p$ non visqueux matérialise
l'effet de la couche limite : faible aux petites incidences, il se
creuse près du bord de fuite (recompression) et en cas de
décollement.

\section{Tracé inférieur : profil et couche limite}

Le tracé du bas affiche le profil \textbf{à l'échelle de la corde}
(repère isométrique : les proportions du profil sont respectées).
Autour de la surface (en blanc) sont tracées les \textbf{frontières de
couche limite} — la surface décalée de l'épaisseur de déplacement
$\delta^*$ le long de la normale extérieure :

\begin{itemize}
    \item extrados en jaune ;
    \item intrados en cyan.
\end{itemize}

L'épaississement de ces frontières vers le bord de fuite traduit la
croissance de la couche limite ; un renflement marqué annonce un
décollement. Le sillage, en aval du bord de fuite, est exclu du tracé.

\section{Encart de coefficients}

Un encart, en haut à droite, rappelle les coefficients intégraux au
point sélectionné : Reynolds, incidence $\alpha$, portance $C_L$,
moment $C_M$, traînée $C_D$ et finesse $C_L/C_D$.

\begin{astuce}
Comparez le $C_p$ visqueux au pointillé non visqueux pour évaluer
d'un coup d'œil l'influence de la viscosité à l'incidence choisie ;
inspectez ensuite les frontières $\delta^*$ pour repérer un
décollement amorcé au bord de fuite.
\end{astuce}

\begin{noteinfo}
Le $C_p$ non visqueux est calculé par une passe XFoil distincte (sans
couche limite). Cette passe est très rapide et peut être désactivée
par le paramètre \code{INVISCID\_CP} (activée par défaut).
\end{noteinfo}
